\documentclass[11pt,a4paper]{article}

% ---- packages -----
\usepackage[a4paper,margin=2.5cm]{geometry}
\usepackage[onehalfspacing]{setspace}
\usepackage[T1]{fontenc}
\usepackage{lmodern}
\usepackage{microtype}
\usepackage{booktabs}
\usepackage{tabularray}
\UseTblrLibrary{booktabs}
\UseTblrLibrary{siunitx}
\usepackage{float}
\usepackage{graphicx}
\usepackage{amsmath, amssymb}
\usepackage{natbib}
\usepackage[hidelinks]{hyperref}
\usepackage{titlesec}
\titlespacing*{\section}{0pt}{1.0\baselineskip}{0.4\baselineskip}
\titlespacing*{\subsection}{0pt}{0.8\baselineskip}{0.3\baselineskip}

% modelsummary helpers
\providecommand{\tinytableTabularrayUnderline}[1]{\underline{#1}}
\providecommand{\tinytableTabularrayStrikeout}[1]{\sout{#1}}

\title{\textbf{Economic Shocks and Civil Conflict in Africa:\\
       A Re-examination of the Miguel-Satyanath-Sergenti\\
       Rainfall-Instrument Approach}}
\author{Alexander Lievens and Tobias Vanoverschelde\\
\small KU Leuven: D0S91a Advanced Applied Econometrics\\
\small Prof.\ Dr.\ Dirk Czarnitzki}
\date{May 2026}

\begin{document}
\maketitle

\begin{abstract}
\noindent
We revisit the seminal instrumental-variables (IV) study of
\citet{MSS2004}, which uses annual rainfall growth as an exogenous
shifter of per-capita GDP growth to identify the causal effect of
economic conditions on civil conflict in 41 sub-Saharan African
countries during 1981 until 1999. Pooled OLS recovers the qualitative
direction of the original association that negative growth coincides
with more conflict, $\hat\beta=-0.51$ ($p=0.02$), but a battery of
modern diagnostics flags the rainfall instruments as weak in our
preferred two-way fixed-effects setting (cluster-robust first-stage
$F=2.89$, well below the Stock-Yogo 10\%-maximum-size threshold of 19.93).
Weak-instrument robust inference \citep{AndersonRubin1949} delivers
a 95\% confidence interval of $[-4.41,\,7.49]$, uninformative about
the sign of the structural parameter. Restoring the country-specific
linear time trends used by MSS in place of year fixed effects almost
doubles the first-stage $F$ to $4.45$ and produces a 2SLS point
estimate of $-1.30$, in line with the OLS sign though still
imprecisely estimated. Most decisively, a dynamic-panel
Arellano-Bond specification \citep{ArellanoBond1991}, which
explicitly accounts for the strong persistence of conflict, yields
a negative and statistically significant coefficient on growth
($-0.29$, $p=0.046$), corroborating the qualitative MSS conclusion
while making clear that the rainfall instrument alone is fragile.
The Wu-Hausman test also fails to reject exogeneity ($p=0.71$),
suggesting that the well-determined OLS estimate may be the most
defensible point estimate in this sample.
\end{abstract}

% ===========================================================================
\section{Introduction}\label{sec:intro}
% ===========================================================================
Whether adverse economic shocks trigger civil conflict is a question
with first-order policy implications for the design of aid, climate
adaptation programmes and post-conflict reconstruction. Ordinary
regressions of conflict on contemporaneous growth confound at least
three sources of endogeneity: reverse causation (conflict destroys
output), measurement error in low-income GDP statistics, and omitted
country traits correlated with both. \citet{MSS2004} address all
three by instrumenting GDP-per-capita growth with year-on-year
rainfall growth, arguing that in agriculture-dependent African
economies rainfall fluctuations are quasi-random shifters of output
that satisfy both relevance and the exclusion restriction.

The MSS finding that a five-percentage-point decline in growth raises
the probability of civil conflict by some twelve percentage
points has been hugely influential, but has also drawn scrutiny.
\citet{Ciccone2011} argues that the link works through rainfall
\emph{levels} rather than growth, and \citet{Sarsons2015} documents
that rainfall predicts conflict downstream of dams that physically
break the agricultural-income channel, casting doubt on the
exclusion restriction.

This paper has two purposes. First, we conduct a careful applied
econometric replication of \citet{MSS2004} that systematically
applies the obtained diagnostics: heteroscedasticity tests, fixed- versus random-effects
choice, weak-instrument tests, overidentification tests,
weak-IV-robust inference, cluster bootstraps, and a control-function
probit. Second, recognising that conflict is highly persistent and
that the rainfall instruments are fragile, we extend the analysis
to a dynamic-panel system-GMM estimator
\citep{ArellanoBond1991,BlundellBond1998} that exploits the panel's
longitudinal structure rather than relying on external instruments.
Our central message is methodological: the MSS conclusion receives
qualified support, but the rainfall instrument is too fragile to
provide compelling causal evidence on its own. The contrast between
specifications is instructive for any applied researcher confronting
a single, marginally relevant instrument.

% ===========================================================================
\section{Conceptual framework}\label{sec:concept}
% ===========================================================================
We adopt the simple ``opportunity-cost'' mechanism of
\citet{Grossman1991} and \citet{ChassangPadroMiquel2009}.
Individuals decide between productive labour, with expected wage $w$
proportional to current income $y$, and predation, with expected
return $\pi$ that is largely independent of $y$. Whenever
$w(y)<\pi$ predation strictly dominates; the conflict-participation
margin is therefore decreasing in $y$. Taking logarithms and first
differences yields the linear approximation
\begin{equation}\label{eq:struct}
   C_{it} = \alpha_i + \delta_t + \beta\, g^{y}_{it}
            + \gamma\,X_{it} + u_{it},
\end{equation}
where $C_{it}=\mathbf{1}\{\text{conflict}_{it}\}$ is a binary
indicator of civil conflict in country $i$ and year $t$,
$g^{y}_{it}=\Delta\ln(y_{it})$ is annual per-capita GDP growth,
$X_{it}$ collects time-varying controls (here lagged
democracy/Polity\,2), and $\alpha_i$, $\delta_t$ are country and
year fixed effects. Theory predicts $\beta<0$.

The empirical challenge is identification of $\beta$: $g^{y}_{it}$
is plausibly endogenous to $u_{it}$. \citet{MSS2004}'s solution is
to instrument $g^{y}_{it}$ with rainfall growth $g^{R}_{it}$ and its
one-period lag. Identification requires (i) \emph{relevance}:
rainfall shifts agricultural output and hence GDP growth, and
(ii) \emph{validity/exclusion?}: rainfall affects conflict only through GDP.
Both assumptions are testable to varying degrees; we report a range
of formal diagnostics below.

% ===========================================================================
\section{Data}\label{sec:data}
% ===========================================================================
We use the replication panel distributed by Edward Miguel through the
Harvard Dataverse (\texttt{doi:10.7910/DVN/27324}), covering 41
sub-Saharan African countries from 1981 to 1999. After constructing
one lag of GDP growth and dropping country-years with missing
controls, our estimation sample is 702 country-year observations
across 41 countries and 18 years (an unbalanced panel). The dependent
variable \texttt{any\_prio} is a PRIO indicator equal to one if a
country experiences at least 25 battle deaths in a given year. The
endogenous regressor is annual GDP-per-capita growth \texttt{gdp\_g};
the excluded instruments are current and lagged rainfall growth
\texttt{GPCP\_g} and \texttt{GPCP\_g\_l}, constructed from the Global
Precipitation Climatology Project. Controls follow MSS: initial log
income, lagged Polity\,2, ethnic and religious fractionalisation, an
oil-exporter dummy, log mountainous terrain, and log lagged
population. Table~\ref{tab:desc} reports descriptive statistics. The
unconditional conflict frequency is 26.9\%, average growth is
slightly negative ($-0.6\%$ p.a.), and the panel is democracy-poor
(mean Polity\,2 of $-3.5$).

\begin{table}[H]
\centering\small
\caption{Descriptive statistics, estimation sample ($N=702$)}\label{tab:desc}
\begin{tblr}{lrrrrr}
\toprule
Variable & $N$ & Mean & SD & Min & Max\\
\midrule
\texttt{any\_prio}  & 702 &  0.269 & 0.444 &  0.000 & 1.000\\
\texttt{gdp\_g}     & 702 & $-0.006$ & 0.066 & $-0.474$ & 0.319\\
\texttt{GPCP\_g}    & 702 &  0.015 & 0.210 & $-0.550$ & 1.680\\
\texttt{y\_0}       & 702 &  1.16  & 0.90  &  0.316 & 4.84\\
\texttt{polity2l}   & 702 & $-3.55$ & 5.55  & $-10$  & 9\\
\texttt{ethfrac}    & 702 &  0.655 & 0.237 &  0.036 & 0.925\\
\texttt{relfrac}    & 702 &  0.487 & 0.185 &  0.000 & 0.783\\
\texttt{Oil}        & 702 &  0.117 & 0.321 &  0     & 1\\
\texttt{lmtnest}    & 702 &  1.58  & 1.43  &  0     & 4.42\\
\texttt{lpopl1}     & 702 &  8.77  & 1.20  &  5.75  & 11.7\\
\bottomrule
\end{tblr}
\end{table}

% ===========================================================================
\section{Econometric specification and methods}\label{sec:methods}
% ===========================================================================
Our reasoning proceeds in five steps that mirror the choices
recommended in \citet{Wooldridge2010}.

\paragraph{(1) Pooled OLS as benchmark.}
We first estimate equation~(\ref{eq:struct}) by Pooled OLS with the
full set of (time-invariant) controls. Because the outcome is binary
this is a Linear Probability Model (LPM) and is heteroscedastic by
construction with $\mathrm{Var}(u_{it})=p_{it}(1-p_{it})$. We test
heteroscedasticity formally with the Breusch--Pagan test and a White
test in the \citet{Wooldridge2010} simplified form, and always report
cluster-robust standard errors clustered at the country level.

\paragraph{(2) Fixed versus random effects.}
The within transformation absorbs the time-invariant controls
(\texttt{y\_0}, \texttt{ethfrac}, \texttt{relfrac}, \texttt{Oil},
\texttt{lmtnest}, \texttt{lpopl1}); we therefore restrict the FE and
RE specifications to time-varying covariates and add year fixed
effects. The presence of country effects is checked with an
$F$-test, and we report the Hausman test of the orthogonality of the
RE-assumed effects to the regressors.

\paragraph{(3) Two-stage least squares.}
We then implement the MSS IV approach with $\{$\texttt{GPCP\_g},
\texttt{GPCP\_g\_l}$\}$ as excluded instruments for \texttt{gdp\_g}.
We report the cluster-robust first-stage $F$ on the excluded
instruments and compare it to the Stock-Yogo critical value (19.93
for 10\% maximum size with two IVs and one endogenous regressor).
We additionally report the Wu-Hausman endogeneity test, the
Sargan/Hansen overidentification test, and the weak-IV-robust
Anderson--Rubin (AR) confidence set, which retains correct coverage
under arbitrarily weak first stages.

\paragraph{(4) Robustness: country trends.}
Because our two-way FE specification produces a markedly weaker first
stage than the original MSS paper, we re-estimate the 2SLS adding
country-specific linear time trends in place of year fixed effects,
the original MSS choice. This sharpens identification by removing
slow country-level secular movements (democratisation, technology
diffusion) without absorbing the rainfall variation.

\paragraph{(5) Dynamic robustness: control function and System-GMM.}
Two further specifications relax the static-binary form. First, we
re-estimate~(\ref{eq:struct}) as a probit with a Rivers-Vuong
control function, plugging in the first-stage residuals as an extra
regressor and reading the $z$-statistic on the residual as a
Hausman-style test of exogeneity \citep{Wooldridge2010}. We prefer
this approach to direct IV-Probit maximum likelihood because, with
41 country dummies and $T=18$, the latter is exposed to
incidental-parameters bias. Second, recognising that conflict is highly
serially correlated, we estimate dynamic specifications including
$C_{i,t-1}$ as a regressor. Within-FE on a dynamic panel suffers
from \citet{Nickell1981} bias of order $O(1/T)$, which with $T=18$
remains non-trivial; we therefore use difference-GMM
\citep{ArellanoBond1991} and system-GMM \citep{BlundellBond1998}
estimators with collapsed instrument matrices and lag depth two
through three to keep the instrument count well below $N=41$
\citep{Roodman2009}. Two-step robust standard errors are reported.
We assess validity with the Sargan test on overidentifying
restrictions and the Arellano-Bond AR(1) and AR(2) tests for
autocorrelation in the differenced residuals.

% ===========================================================================
\section{Results}\label{sec:results}
% ===========================================================================
\subsection{Main estimates}
Table~\ref{tab:main} collects the seven static specifications.
Pooled-OLS column (1) reproduces the headline MSS direction: a one
standard-deviation drop in growth ($\approx 6.6$ percentage points)
raises the conflict probability by roughly $0.51\times 0.066\approx
3.4$ percentage points, statistically significant at the 5\% level.
Two-way fixed-effects column (2) and random-effects column (3)
shrink the coefficient towards zero ($\hat\beta=-0.46$, $-0.44$)
with slightly widened standard errors, as expected when we strip out
between-country variation. The Hausman test fails to reject RE
($p=0.998$); but the $F$-test for country effects is large
($F=14.9$, $p<0.001$); given the strong rejection of pooling,
we maintain FE as the preferred estimator. The reduced-form column (4)
shows that rainfall growth enters the conflict equation only weakly
(lagged rainfall has $p=0.08$); the first stage column (5) confirms
that current-year rainfall growth has a statistically significant
but quantitatively modest effect on GDP growth.

Column (6), our baseline 2SLS, is the most striking result. The
point estimate on \texttt{gdp\_g} is $+0.14$ with a clustered
standard error of $1.33$: the wrong sign, no statistical
significance, and a confidence interval covering essentially all
economically relevant magnitudes. This collapse of the IV estimate
relative to the well-determined POLS estimate is the canonical
fingerprint of a weak instrument. Column (7) re-estimates the same
model with country FE plus country-specific linear time trends---the
MSS specification---and yields $\hat\beta=-1.30$ (SE $1.07$,
$p=0.23$): the correct (negative) sign, larger in magnitude than
OLS, but still imprecisely estimated. The qualitative direction of
the original MSS finding is therefore highly sensitive to whether
secular country-level trends are modelled.

% Table inserted from the analysis script
\begin{table}[H]
\centering
\small
\input{tables/main_results.tex}
\caption{Effect of GDP growth on civil conflict in Africa, 1981--1999. Models~(1)--(6) include year fixed effects (where applicable); model~(7) replaces year FE with country-specific linear time trends. Cluster-robust SE in parentheses.}
\label{tab:main}
\end{table}

\subsection{Specification diagnostics}
Table~\ref{tab:diag} reports the formal diagnostics, grouped by what
they test.

\emph{Heteroscedasticity.}
The Breusch--Pagan test overwhelmingly rejects ($\chi^2=142.5$,
$p<10^{-16}$), confirming the necessity of robust standard errors in
the LPM.

\emph{Panel structure.}
The $F$-test of country fixed effects strongly rejects the null
that they are jointly zero ($F=14.9$, $p<10^{-16}$); the Hausman
test does not distinguish FE from RE ($p=0.998$). The classical
Hausman statistic is, however, unreliable in this setting: its
construction assumes that the FE estimator is efficient under
$H_0$, which fails once standard errors are clustered, and the test
also has limited power in moderately unbalanced panels. We
therefore default to FE following \citet{Wooldridge2010}.

\emph{Weak instruments.}
The cluster-robust first-stage $F$ on the two excluded instruments
is just $2.89$ in our baseline 2-way FE specification, far below
the \citet{StaigerStock1997} rule-of-thumb of 10 or the
Stock--Yogo \citep{StockYogo2005} 10\%-maximum-size threshold of
19.93 (the corresponding 10\%-relative-bias critical value is
$\approx 11$, which the baseline $F$ also fails to meet). Switching to the MSS-style country-trend specification
lifts the $F$ to $4.45$ (cluster-robust Wald) or $7.48$
(homoscedastic $F$), notably stronger but still below conventional
thresholds.

\emph{Endogeneity.}
The Wu-Hausman exogeneity test fails to reject ($\chi^2=0.14$,
$p=0.71$); the control-function probit residual $\hat v$ is
statistically insignificant ($z=-0.73$, $p=0.46$). In this sample
we have no positive evidence that OLS is inconsistent.

\emph{Overidentification.}
The Sargan test for the two-instrument design does not reject
($\chi^2=2.40$, $p=0.12$), but with only one overidentifying
restriction this test has low power and provides only weak
reassurance about the exclusion restriction.

\emph{Weak-IV-robust inference.}
The Anderson-Rubin 95\% confidence set for $\beta$ is
$[-4.41,\,7.49]$, essentially uninformative about either the sign or
the magnitude of the structural parameter. A non-parametric cluster
bootstrap (500 replications, resampling countries with replacement)
yields a closely comparable percentile interval of $[-3.52,\,3.51]$.

\begin{table}[H]
\centering\small
\caption{Specification and identification diagnostics}\label{tab:diag}
\begin{tblr}{lrr}
\toprule
Statistic & Value & $p$-value \\
\midrule
\SetCell[c=3]{l}\textit{Heteroscedasticity} & & \\
Breusch--Pagan                                     & $142.52$ & $<10^{-16}$\\
\midrule
\SetCell[c=3]{l}\textit{Panel structure} & & \\
$F$-test country FE (H$_0$: $c_i=0$)               & $14.90$  & $<10^{-16}$\\
Hausman FE vs RE                                   &  $0.04$  & $0.998$ \\
\midrule
\SetCell[c=3]{l}\textit{Weak instruments} & & \\
First-stage $F$ baseline (cluster-robust)          &  $2.89$  & $0.056$ \\
First-stage $F$ MSS-trends (cluster-robust Wald)   &  $4.45$  & $0.012$ \\
First-stage $F$ MSS-trends (homoscedastic)         &  $7.48$  & $0.0006$ \\
Stock--Yogo 10\%-maximum-size critical value      & $19.93$  & --- \\
\midrule
\SetCell[c=3]{l}\textit{Endogeneity} & & \\
Wu--Hausman (AER \texttt{ivreg})                   &  $0.14$  & $0.705$ \\
Control-function probit $z(\hat v)$                & $-0.73$  & $0.464$ \\
\midrule
\SetCell[c=3]{l}\textit{Overidentification} & & \\
Sargan (AER \texttt{ivreg})                        &  $2.40$  & $0.121$ \\
\midrule
\SetCell[c=3]{l}\textit{Weak-IV-robust inference} & & \\
Anderson--Rubin 95\% CI for $\beta$                & $[-4.41,\,7.49]$ & --- \\
Cluster-bootstrap 95\% CI for $\beta^{IV}$         & $[-3.52,\,3.51]$ & --- \\
\bottomrule
\end{tblr}
\end{table}

\subsection{Robustness: dynamic panel-data estimators}
The dynamic-panel results in Table~\ref{tab:gmm} are the most
encouraging for the substantive MSS conclusion. The lagged conflict
indicator carries a large and highly significant coefficient
($0.63$, $p<0.001$ in system-GMM), confirming substantial conflict
persistence and indicating that the static specifications are
mis-specified. Conditional on this dynamic structure, the
Arellano-Bond estimator delivers $\hat\beta=-0.29$ with $p=0.046$;
the Blundell-Bond system-GMM yields a similar point estimate
$\hat\beta=-0.30$ but no longer statistically significant
($p=0.19$). The Sargan tests do not reject ($p=0.51$ and $0.09$
respectively); AR(2) tests do not detect second-order
autocorrelation in differences ($p=0.50$ and $0.47$); both
validations support the identifying restrictions. The qualitative
MSS result: a negative relationship between growth and
conflict, thus reappears once the dynamic structure of conflict is
properly modelled, even without an external instrument.

\begin{table}[H]
\centering
\small
\input{tables/gmm_results.tex}
\caption{Dynamic panel-data estimators of conflict on growth. Two-step robust standard errors. Collapsed instrument matrix with lag depth 2--3.}
\label{tab:gmm}
\end{table}

% ===========================================================================
\section{Discussion and conclusion}\label{sec:conc}
% ===========================================================================
This essay illustrates a tension that recurs in applied
microeconometrics: a story that is compelling, theoretically well
motivated and empirically detectable in OLS does not always survive
strict instrumental-variables identification. Four lessons stand
out.

First, the rainfall instruments are markedly weaker in our two-way
FE specification ($F=2.89$) than under the MSS country-trends design
($F=4.45$). This is a useful reminder that the strength of an
instrument is specification-dependent, and that seemingly small
modelling choices (year fixed effects versus country trends)
materially change identification.

Second, when an instrument is weak, conventional 2SLS $t$-statistics
are unreliable and the only valid inference is weak-IV-robust
\citep{AndersonRubin1949,StaigerStock1997}. Reporting only the
2SLS point estimate would have been misleading: our baseline 2SLS
estimate is positive while every other reasonable specification
points negative. The Anderson-Rubin confidence set, which makes no
strength-of-identification assumption, is consequently very wide.

Third, the Wu-Hausman test and the control-function probit residual
both fail to reject exogeneity of GDP growth. Combined with the
weak first stage, this would justify reporting the well-determined
POLS estimate of $-0.51$ as the most defensible point estimate in
this sample, treating the IV strategy as an inconclusive robustness
check rather than the primary identification.

Fourth, leveraging the panel dimension via dynamic-panel GMM offers
an identification route that does not depend on an external
instrument at all. Conditional on lagged conflict, both
Arellano-Bond and Blundell-Bond recover negative growth
coefficients near $-0.30$, the first significant at the 5\% level.
Our overall reading is that MSS's substantive conclusion, namely that adverse
economic shocks raise the probability of civil conflict in
Africa, receives qualified support, but the rainfall instrument is
too fragile to provide compelling causal evidence on its own.

Limitations of our analysis include the binary measure of conflict
(LPM versus probit/logit), the focus on contemporaneous shocks
(richer distributed lags are plausible), and the omission of
conflict-spillover spatial dynamics emphasised in subsequent
literature. Each of these would be natural extensions for a
longer-form study.

% ---------------------- bibliography ----------------------
\renewcommand{\refname}{References}
\begin{thebibliography}{99}\small

\bibitem[Anderson and Rubin(1949)]{AndersonRubin1949}
Anderson, T.W., and H.\ Rubin (1949). ``Estimation of the parameters
  of a single equation in a complete system of stochastic equations.''
  \emph{Annals of Mathematical Statistics} 20(1): 46--63.

\bibitem[Arellano and Bond(1991)]{ArellanoBond1991}
Arellano, M., and S.\ Bond (1991). ``Some tests of specification for
  panel data: Monte Carlo evidence and an application to employment
  equations.'' \emph{Review of Economic Studies} 58(2): 277--297.

\bibitem[Blundell and Bond(1998)]{BlundellBond1998}
Blundell, R., and S.\ Bond (1998). ``Initial conditions and moment
  restrictions in dynamic panel data models.''
  \emph{Journal of Econometrics} 87(1): 115--143.

\bibitem[Chassang and Padr\'o i Miquel(2009)]{ChassangPadroMiquel2009}
Chassang, S., and G.\ Padr\'o i Miquel (2009). ``Economic shocks and
  civil war.'' \emph{Quarterly Journal of Political Science} 4(3):
  211--228.

\bibitem[Ciccone(2011)]{Ciccone2011}
Ciccone, A.\ (2011). ``Economic shocks and civil conflict: A
  comment.'' \emph{American Economic Journal: Applied Economics}
  3(4): 215--227.

\bibitem[Grossman(1991)]{Grossman1991}
Grossman, H.I.\ (1991). ``A general equilibrium model of
  insurrections.'' \emph{American Economic Review} 81(4): 912--921.

\bibitem[Miguel et al.(2004)]{MSS2004}
Miguel, E., S.\ Satyanath, and E.\ Sergenti (2004). ``Economic
  shocks and civil conflict: An instrumental variables approach.''
  \emph{Journal of Political Economy} 112(4): 725--753.

\bibitem[Nickell(1981)]{Nickell1981}
Nickell, S.\ (1981). ``Biases in dynamic models with fixed
  effects.'' \emph{Econometrica} 49(6): 1417--1426.

\bibitem[Roodman(2009)]{Roodman2009}
Roodman, D.\ (2009). ``A note on the theme of too many instruments.''
  \emph{Oxford Bulletin of Economics and Statistics} 71(1): 135--158.

\bibitem[Sarsons(2015)]{Sarsons2015}
Sarsons, H.\ (2015). ``Rainfall and conflict: A cautionary tale.''
  \emph{Journal of Development Economics} 115: 62--72.

\bibitem[Staiger and Stock(1997)]{StaigerStock1997}
Staiger, D., and J.H.\ Stock (1997). ``Instrumental variables
  regression with weak instruments.'' \emph{Econometrica} 65(3):
  557--586.

\bibitem[Stock and Yogo(2005)]{StockYogo2005}
Stock, J.H., and M.\ Yogo (2005). ``Testing for weak instruments in
  linear IV regression.'' In: \emph{Identification and Inference for
  Econometric Models: Essays in Honor of Thomas Rothenberg}, ed.\
  D.W.K.\ Andrews and J.H.\ Stock, Cambridge University Press,
  80--108.

\bibitem[Wooldridge(2010)]{Wooldridge2010}
Wooldridge, J.M.\ (2010). \emph{Econometric Analysis of Cross
  Section and Panel Data.} 2nd ed., MIT Press.

\end{thebibliography}

\end{document}
